\chapter{Perimenopause}

\medskip
\noindent\textit{\textbf{Under review.} This chapter is an AI-assisted educational draft; it is not a substitute for professional medical advice.}

\section*{Definition and Overview}
Perimenopause concerns reproductive, sexual, menstrual, pregnancy, or menopause-related health. Symptoms and management vary with age, pregnancy status, underlying conditions, and personal goals.

\section*{Clinical Features}
Possible features include changes in bleeding, pain, discharge, fertility, hormonal symptoms, pregnancy symptoms, or sexual function.

\section*{When to Seek Medical Care}
Seek medical review for new, persistent, worsening, or function-limiting symptoms. Call 000 for severe breathing difficulty, collapse, sudden neurological change, severe chest or abdominal pain, or any rapidly worsening illness.

\section*{Diagnosis and Investigations}
Assessment may include pregnancy testing, examination, blood tests, sexual-health testing, ultrasound, and other imaging when indicated.

\section*{Management}
Management is individualised and may include education, symptom treatment, contraception or fertility care, infection treatment, and referral to a GP or reproductive-health specialist.

\section*{Complications}
Complications depend on the underlying cause and may include progression, effects on other organs, treatment adverse effects, or reduced quality of life.

\section*{Prognosis and Self-Care}
Outcome depends on the cause, severity, complications, and response to treatment. Follow the care plan, keep appointments, avoid known triggers, and seek help if symptoms worsen.

\section*{Expanded clinical context}
This chapter addresses perimenopause, a topic that should be interpreted in the context of age, baseline health, medicines, comorbidities, goals, and access to care. It supports discussion with a qualified clinician and does not establish a diagnosis.

\section*{Assessment and decision points}
Review onset, duration, triggers, relevant exposures, physical findings, associated symptoms, functional impact, and immediate safety. Record the main concern, time course, risk factors, red flags, and the person’s questions. Use targeted investigations when their result is likely to change management.

\section*{Management and follow-up}
Care may include education, practical support, treatment of contributing conditions, medication or a procedure when indicated, and a shared follow-up plan. Explain expected improvement, possible adverse effects, and when reassessment is needed. Avoid starting, stopping, or sharing prescription medicines without professional advice.

\section*{Safety-netting}
Seek urgent help for severe or rapidly worsening symptoms, collapse, breathing difficulty, severe pain, confusion, dehydration, uncontrolled bleeding, or any immediate danger.

\section*{Practical prevention}
Use plain-language information, supportive routines, risk reduction, and professional review when symptoms persist, recur, or impair daily life. Keep written instructions and confirm how to access help outside routine appointments.

\section*{Limitations}
This educational expansion should be checked against current local guidance and authoritative topic-specific sources.
\clearpage
